\documentclass[10pt,a4paper]{amsart}
\usepackage[margin=19mm]{geometry}
\usepackage{amsmath,amssymb,mathtools}
\usepackage[scr=rsfs,cal=euler]{mathalfa}
\usepackage{xfrac}
\usepackage[unicode,hidelinks,bookmarks=false]{hyperref}
\newcommand{\ie}{i.e.\ }
\newcommand{\cf}{cf.\ }
\newcommand{\resp}{respectively\ }
\newcommand{\Subsection}{\subsection}
\providecommand{\nobreakdash}{\nobreak\hbox{-}}
\setlength{\emergencystretch}{2em}
\multlinegap0pt
\usepackage{comment}
\usepackage[all]{xy}
\theoremstyle{plain}
\newtheorem{thm}{Theorem}
\newtheorem{prop}[thm]{Proposition}
\newtheorem{lem}[thm]{Lemma}
\newtheorem{cor}[thm]{Corollary}
\newtheorem{conj}[thm]{Conjecture}
\theoremstyle{definition}
\newtheorem{defn}[thm]{Definition}
\theoremstyle{remark}
\newtheorem{remark}[thm]{Remark}
\newcommand{\Q}{\mathbf{Q}}
\newcommand{\Z}{\mathbf{Z}}
\def\cL{{\mathcal L}}
\def\cM{\mathcal M}
\def\cO{\mathcal O}
\DeclareMathOperator{\la}{\lambda}
\title[Anticyclotomic Euler systems]{Anticyclotomic Euler system over biquadratic fields: the ring-class exact sequence attached to a squarefree ideal of the biquadratic field}
\author[K. T. Do]{Kim Tuan Do}
\date{Source: arXiv:2409.19819v2 (CC BY 4.0)}
\begin{document}
\maketitle
\noindent\emph{Source and licence.} Anticyclotomic Euler system over biquadratic fields. Version arXiv:2409.19819v2 (CC BY 4.0), distributed under the Creative Commons Attribution 4.0 International licence, \url{http://creativecommons.org/licenses/by/4.0/}, as recorded in the arXiv licence metadata for this exact version and shown on the abstract page \url{https://arxiv.org/abs/2409.19819v2}. Full licence notice: licenses/arxiv-2409.19819-CC-BY-4.0.txt.\par\smallskip

\noindent\textit{Editorial scope.} Counted unit: the article's Proposition with source label \texttt{ringclass} (source0.tex lines 797--817), the exact sequence that decomposes the p-part of the ray class group of the biquadratic field modulo a squarefree ideal into the two partial ray class groups and the class group modulo that ideal, together with the description of the projection composed with the reciprocity map by a Frobenius element; delivered with the article's complete original proof of it (source0.tex lines 819--837, one \texttt{proof} environment of 19 lines, adjacent to the statement). No mathematics is written, repaired or re-derived: statement and proof are the article's own text line for line, in the article's own notation (the ray class groups and their p-parts, the reciprocity map, the ideal and its norm, the Frobenius element), and no retained line is substituted, re-flowed or omitted. Required context retained uncounted: the labelled display of the decomposition of the two partial ray class groups (lines 788--790), which the counted proof invokes by label at its last step. Nothing else of the article is retained: its main theorems, its Euler-system construction and its other propositions are omitted. No citation command occurs in any retained line, so the article's bibliography is not delivered and the wrapper carries no bibliography entry, although the article ships one. Every cross-reference inside the retained spans resolves inside the delivered document. Editorial formatting only: an AMS \texttt{amsart} preamble that restates the environments and macros the retained lines use, namely the article's declarations of Theorem, Proposition, Lemma, Corollary, Conjecture, Definition and Remark sharing one counter, the article's two macros for the rational and the integer field, its three macros for the calligraphic letters, its operator for the prime element, the \texttt{xy} package that draws the article's commutative diagrams, and the \texttt{comment} package that the article itself uses and that suppresses one commented-out diagram carried verbatim inside the counted statement. The article's own source declares the class \texttt{amsart} as well; no class or style file is shipped with this package. Numbering differs from the article, because the article resets its shared theorem counter per subsection and because only a subset of environments is retained: the counted proposition prints as Proposition 1. No picture environment and no graphical inclusion occurs in any retained line, and the package delivers no image asset; the commutative diagrams of the retained spans are drawn by the retained diagram code itself. One bundled source file, \texttt{sources/166008/source0.tex}, accompanies the delivered item as the exact source of every retained span. Every retained span is recorded in \texttt{delivery-evidence.json} with method \texttt{exact\_tex}.
\begin{equation}\label{eq:H-m-decompose}
    ^{1}H_{\mu_1}^{(p)}\times {^{2}H}_{\mu_2}^{(p)} \xrightarrow{\simeq} (\cO_{K_0}/\mu_3\cO_{K_0})^{\times,(p)} \twoheadrightarrow {^{0}H_{\mu_3}^{(p)}}.
\end{equation}

\begin{prop}\label{ringclass} Suppose $p\nmid 6h_{K_0}$ and $\mu_3$ is a squarefree ideal of $\cO_{K_3}$ of norm $m$, where $m$ is divisible only by primes that are split in $K_0$. Using (\ref{eq:H-m-decompose}), we have the following short exact sequence:
\begin{comment}
    \[
\xymatrix{
(\cO_{K_0}/\mu_3\cO_{K_0})^{\times,(p)} \ar[r]^{\simeq}&{^{1}H_{\mu_1}^{(p)}}\times {^{2}H_{\mu_2}^{(p)}}\ar@{->>}[d]^{w}& & \\
(\cO_{K_3}/\mu_3\cO_{K_3})^{\times, (p)}\ar[r]\ar[u]^{\Delta}&{^{0}H_{\mu_3}^{(p)}}\ar[r]^{e}&H[\mu_3]^{(p)}\ar[r]&1.
}
\]
\end{comment}
\[
1\rightarrow(\cO_{K_3}/\mu_3\cO_{K_3})^{\times, (p)}\xrightarrow{\Delta}{^{1}H}_{\mu_1}^{(p)}\times {^{2}H}_{\mu_2}^{(p)}\xrightarrow{e\circ w} H[\mu_3]^{(p)}\rightarrow 1.
\]
Here, the map $\Delta$ uses the identifications 
\[(\cO_{K_3}/\mu_3\cO_{K_3})^{\times,(p)} \simeq (\cO_{K_0}/\cM_4\cO_{K_0})^{\times,(p)},\qquad (\cO_{K_3}/\mu_3\cO_{K_3})^{\times,(p)} \simeq (\cO_{K_0}/\cM_3\cO_{K_0})^{\times,(p)}.
\]
Moreover, if $(\ell)=\cL_1\cL_2\cL_3\cL_4$ is a prime that splits in $K_0$ and is coprime to $m$, the projection $e\circ w$ (defined in the proof below) sends 
\begin{align*}
      \quad  [\la_1]\times[\bar{\la}_2]&\mapsto \mathrm{Frob}_{\cL_4/\la_3} 
    \end{align*}
where $\mathrm{Frob}_{\cL_4/\la_3} $ is the Frobenius element of $\cL_4$ in $H[\mu_3]^{(p)}$. 
\end{prop}

\begin{proof}
We have the following exact diagram: 
\[
\xymatrix{
\cO_{K_0}^\times\ar[r]\ar[d]&(\cO_{K_0}/\mu_3\cO_{K_0})^\times\ar[r]\ar[d]&^{0}H_{\mu_3}\ar[r]\ar[d]&^{0}H_1\ar[r]\ar[d]&1\\
\cO_{K_0}^\times/\cO_{K_3}^\times\ar[r]&(\cO_{K_0}/\mu_3\cO_{K_0})^\times/(\cO_{K_3}/\mu_3\cO_{K_3})^\times\ar[r]&H[\mu_3]\ar[r]&^{0}H_1\ar[r]&1.
}
\]
Taking the $p-$part of this, using the assumption that $p\nmid 6h_{K_0}$ together with the fact that $|\cO_{K_0}^\times/\cO_{K_3}^\times|$ is a power of $2$, we obtain the following diagram:
\[
\xymatrix{
\cO_{K_0}^{\times}\otimes \frac{\Q_p}{\Z_p}\ar[r]\ar[d]&(\cO_{K_0}/\mu_3\cO_{K_0})^{\times,(p)}\ar[r]^{w}\ar[d]&^{0}H_{\mu_3}^{(p)}\ar[r]\ar[d]^{e}&1\\
1\ar[r]&(\cO_{K_0}/\mu_3\cO_{K_0})^{\times,(p)}/(\cO_{K_3}/\mu_3\cO_{K_3})^{\times,(p)}\ar[r]&H[\mu_3]^{(p)}\ar[r]&1.
}
\]
Using the middle arrow and the identification \ref{eq:H-m-decompose}, we can show the first exact sequence.

One can show the second part by noting that $\cL_4$ is a prime of $K_0$ lying above both $\la_1$ (a prime of $K_1$) and $\bar{\la}_2$ (a prime of $K_2$).
\end{proof}

\end{document}
